\documentclass[a4paper]{article}
\usepackage[margin=3cm]{geometry}
\usepackage{graphicx}
\usepackage{qrcode}

\usepackage{hyperref}
\hypersetup{
	colorlinks=true,
	linkcolor=blue,
	filecolor=magenta,
	urlcolor=cyan,
	hyperindex=true,
	linktocpage=true,%makes the page number be the link in the index
}
\usepackage{listings}
\lstset{
	basicstyle=\tiny\ttfamily,
	columns=fullflexible,
	keepspaces=true,
	breaklines=true,
	frame=single,
	showstringspaces=true,
	tabsize=4,
	showspaces=true,
	showtabs=true,
}

\renewcommand{\arraystretch}{1.2} %adds a bit of margin to table cells
\renewcommand{\thefootnote}{[\arabic{footnote}]} %makes footnotes look wikipedia-style

\title{Finespun standard environment}

\author{Chaya2766}

\date{last updated \today}

\begin{document}
	\sffamily
	\hyphenchar\font=-1
	\sloppy
	
	\maketitle
	
	\begin{center}
		\qrcode[height=2cm]{https://chaya2766.github.io/stareater-expanse-wiki/}
	\end{center}
	
	\textbf{
		\sloppy
		\hyphenchar\font=-1
		}
	
	\vfill
	
	\begin{abstract}
		Finespun settlements opt to use a pure oxygen to deliver similar oxygen availability as terran sea level atmosphere.
		
		\medskip
		
		The environment is specified as 0.21 bar of pure oxygen and ambient temperature of 290 kelvin.
		
		\medskip
		
		This document goes over some of the details of this environment and its possible effects on people and daily activities.
	\end{abstract}
	
	\vfill
	
	\begin{quote}
		\begin{center}
			\textbf{DISCLAIMER AND WARNING}
		\end{center}
		
		This is not a scientific paper or engineering advice.
		This document is only intended as creative world‑building and has not been peer‑reviewed or approved for any practical application.
		Use of these ideas in real‑world designs or other serious application is entirely at the reader’s discretion and risk.
	\end{quote}
	
	%\noindent \rule{\linewidth}{1pt}
	\vfill
	
	\tableofcontents
	
	\pagebreak
	
	\section{Water boiling point (\& custom temperature units)}
	
	The following comes from the "Clausius-Clapeyron relation" which can provide the boiling point of a substance to any pressure given a known boiling point at a known pressure.
	
	$$ \ln\left(\frac{P_1}{P_0}\right) = \frac{L}{R \cdot (\frac{1}{T_0} - \frac{1}{T_1}) } $$
	
	$$ T_1 = \frac{1}{\frac{1}{T_0} - \ln\left(\frac{P_1}{P_0}\right) \cdot \frac{R}{L}} $$
	
	where:
	
	\begin{itemize}
		\item $T_0$ is the original observed boiling point
		
		\item $P_0$ is the original pressure at which the original boiling point occurred
		
		\item $P_1$ is the desired pressure
		
		\item $T_1$ is the boiling point of the substance at the desired pressure
		
		\item $L$ is the molar enthalpy of vaporization of the substance
		
		\item $R$ is the ideal gas constant ($\approx$8.314462618 J/(K*mol))
	\end{itemize}
	
	\begin{center}
		For water the known values are as follows:
		
		\begin{tabular}{c | c | c }
			$T_0$ & $P_0$ & $L$ \\\hline
			373.15 K & 101 325 Pa & 40 660 J/mol
		\end{tabular}
	\end{center}
	
	The desired environment has a pressure of 0.21 bar (21kPa), which will be the $P_1$ to plug into the equation:
	
	$$ T_1 = \frac{1}{\frac{1}{373.15 K} - \ln\left(\frac{21 000 pa}{101 325 Pa}\right) \cdot \frac{8.314462618 J/(K*mol)}{40 660 J/mol}} = 333.1432084 K$$
	
	Finespun's 0.21 bar atmosphere causes water to boil at 333.143K $\approx$ 59.993$^o$C $\approx$ 139.988$^o$F
	
	\bigskip
	
	\begin{center}
		\rule{0.8\linewidth}{1pt}
		
		\textbf{custom temperature unit}
	\end{center}
	
	A new temperature should be defined which would be convenient for everyday use, Finespun's equivalent of celcius which sets 0 at the freezing point and 100 at the boiling point.
	It will be defined in terms of degrees kelvin.
	
	$$ (273.15K - a) \cdot b = 0^oW , (333.143K - a) \cdot b = 10^oW $$
	
	To make these equations work, a = 273.15K, b = 1.666861134 $\frac{^oW}{K}$.
	
	As such, any temperature expressed in kelvin can be converted into this new custom unit as such:
	
	$$ w = (k - 273.15K) \cdot 1.666861134 \frac{^oW}{K} $$
	
	where:
	
	\begin{itemize}
		\item $w$ is the temperature in $^oW$
		
		\item $k$ is the temperature in degrees kelvin
	\end{itemize}
	
	\bigskip
    
    From $^oC$ to $^oW$ the conversion is closely approximated by dividing by 0.6$\frac{^oC}{^oW}$
    
	\pagebreak
	
	\section{todo}
	
\end{document}